\documentclass[a4paper,12pt]{ctexart}
\usepackage{xcolor}
\usepackage[top=1.8cm, bottom=1.8cm, left=1.8cm, right=1.8cm]{geometry}
\usepackage{array}
\usepackage{listings}
\usepackage{hyperref}
\hypersetup{colorlinks=true,linkcolor=blue!70!black}

\lstset{
  basicstyle=\ttfamily\small,
  keywordstyle=\color{blue},
  commentstyle=\color{gray},
  stringstyle=\color{red},
  showstringspaces=false,
  frame=single,
  numbers=left,
  numberstyle=\tiny,
  breaklines=true
}

\title{\textcolor{blue!70!black}{\Huge\textbf{数据结构与算法初步}}}
\author{Python 进阶——像计算机科学家一样思考}
\date{}

\begin{document}
\maketitle
\thispagestyle{empty}

\section*{写在前面}

你已经学会了 Python 的基本语法，能写 100 行左右的程序了。现在，我们要进入计算机科学最核心的领域：\textbf{数据结构与算法}。

简单来说：
\begin{itemize}
  \item \textbf{数据结构} —— 数据在计算机里怎么\textbf{组织}（怎么存、怎么找）
  \item \textbf{算法} —— 解决问题的方法和\textbf{步骤}（怎么算、怎么处理）
\end{itemize}

这两样东西决定了你的程序是快还是慢，是优雅还是笨拙。

\section*{第一课：时间复杂度——程序跑多快？}

两个程序都能做同一件事，但一个用了 0.1 秒，另一个用了 10 秒。为什么呢？因为它们的\textbf{算法效率}不同。

\subsection*{大 O 表示法}

计算机科学家用"大 O"来表示算法的速度。它描述的是：\textbf{输入变大时，运行时间增长得多快}。

\begin{center}
\begin{tabular}{|c|c|c|c|}
\hline
\textbf{符号} & \textbf{名称} & \textbf{含义} & \textbf{举例} \\ \hline
$O(1)$ & 常数时间 & 不管输入多大，时间不变 & 访问列表第一个元素 \\ \hline
$O(\log n)$ & 对数时间 & 输入翻倍，只多一步 & 二分查找 \\ \hline
$O(n)$ & 线性时间 & 输入多大，时间就多长 & 遍历列表 \\ \hline
$O(n \log n)$ & 线性对数 & 比线性稍慢一点 & 归并排序、快排 \\ \hline
$O(n^2)$ & 平方时间 & 输入翻倍，时间变 4 倍 & 冒泡排序 \\ \hline
$O(2^n)$ & 指数时间 & 输入加 1，时间翻倍 & 斐波那契朴素递归 \\ \hline
\end{tabular}
\end{center}

\subsection*{直观感受}

假设你的计算机每秒能执行 100 万次操作：

\begin{center}
\begin{tabular}{|c|c|c|c|}
\hline
\textbf{复杂度} & \textbf{$n=10$} & \textbf{$n=100$} & \textbf{$n=1000$} \\ \hline
$O(1)$ & 0.000001 秒 & 0.000001 秒 & 0.000001 秒 \\ \hline
$O(\log n)$ & 0.000003 秒 & 0.000007 秒 & 0.00001 秒 \\ \hline
$O(n)$ & 0.00001 秒 & 0.0001 秒 & 0.001 秒 \\ \hline
$O(n \log n)$ & 0.00003 秒 & 0.0007 秒 & 0.01 秒 \\ \hline
$O(n^2)$ & 0.0001 秒 & 0.01 秒 & 1 秒 \\ \hline
$O(2^n)$ & 0.001 秒 & $4 \times 10^{16}$ 年 & - \\ \hline
\end{tabular}
\end{center}

\textbf{关键：} 选对算法比买更好的电脑重要得多！

\subsection*{练习：判断复杂度}

以下代码的时间复杂度是多少？

\begin{lstlisting}[language=Python]
# 程序 A
def sum_list(nums):
    total = 0
    for n in nums:
        total += n
    return total

# 程序 B
def print_pairs(nums):
    for i in nums:
        for j in nums:
            print(i, j)

# 程序 C
def binary_search(lst, target):
    low, high = 0, len(lst) - 1
    while low <= high:
        mid = (low + high) // 2
        if lst[mid] == target:
            return mid
        elif lst[mid] < target:
            low = mid + 1
        else:
            high = mid - 1
    return -1
\end{lstlisting}

答案：A 是 $O(n)$，B 是 $O(n^2)$，C 是 $O(\log n)$（假设列表已排序）。

\section*{第二课：数组 vs 链表}

\subsection*{数组（Python 列表）}

你已经在 Python 中大量使用了列表（其实是动态数组）。

\textbf{特点：}
\begin{itemize}
  \item 在内存中是\textbf{连续}的一整块
  \item 通过索引访问非常快：$O(1)$
  \item 在中间插入/删除慢：需要移动后面的元素，$O(n)$
  \item 尾部添加快：$O(1)$
\end{itemize}

\subsection*{链表}

链表像\textbf{一列火车}：每节车厢（节点）里装着数据，还有一根链条连到下一节。

\begin{lstlisting}[language=Python]
class Node:
    def __init__(self, value):
        self.value = value
        self.next = None

class LinkedList:
    def __init__(self):
        self.head = None
    
    def append(self, value):
        if not self.head:
            self.head = Node(value)
            return
        current = self.head
        while current.next:
            current = current.next
        current.next = Node(value)
    
    def print_all(self):
        current = self.head
        while current:
            print(current.value, end=" -> ")
            current = current.next
        print("None")
\end{lstlisting}

\textbf{特点：}
\begin{itemize}
  \item 在内存中\textbf{不连续}，分散在各处
  \item 通过索引访问慢：必须从头遍历，$O(n)$
  \item 任意位置插入/删除快：只需要改链条，$O(1)$
\end{itemize}

\subsection*{对比}

\begin{center}
\begin{tabular}{|c|c|c|}
\hline
\textbf{操作} & \textbf{数组} & \textbf{链表} \\ \hline
按索引访问 & $O(1)$ & $O(n)$ \\ \hline
尾部插入 & $O(1)$ & $O(1)$（有尾指针） \\ \hline
头部插入 & $O(n)$ & $O(1)$ \\ \hline
中间插入 & $O(n)$ & $O(1)$（已知位置） \\ \hline
\end{tabular}
\end{center}

\textbf{记住了：} 没有"最好的"数据结构，只有"最适合当前场景"的数据结构。

\section*{第三课：栈与队列}

\subsection*{栈——后进先出}

\textbf{想象一叠盘子：} 你总是拿最上面的盘子（最后放上去的），放也放在最上面。

Python 列表可以当栈用：

\begin{lstlisting}[language=Python]
stack = []
stack.append("第一本")   # 入栈
stack.append("第二本")
stack.append("第三本")

print(stack.pop())       # 出栈：第三本
print(stack.pop())       # 第二本
print(stack.pop())       # 第一本
\end{lstlisting}

\textbf{生活中的栈：} 浏览器的"后退"按钮（你点进的每个页面都被压入栈中，点后退就弹出上一页）、撤销操作（Ctrl+Z）。

\subsection*{队列——先进先出}

\textbf{想象排队买奶茶：} 先来的人先买到，后来的人排后面。

\begin{lstlisting}[language=Python]
from collections import deque

queue = deque()
queue.append("小明")   # 入队
queue.append("小红")
queue.append("小刚")

print(queue.popleft())  # 出队：小明（最先来的）
print(queue.popleft())  # 小红
print(queue.popleft())  # 小刚
\end{lstlisting}

\textbf{生活中的队列：} 打印机任务队列、游戏匹配排队。

\subsection*{练习：括号匹配}

用栈检查一个字符串中的括号是否匹配：\texttt{"()[]\{\}"} 是匹配的，\texttt{"([)]"} 是不匹配的。

\begin{lstlisting}[language=Python]
def is_valid(s):
    stack = []
    pairs = {")": "(", "]": "[", "}": "{"}
    
    for ch in s:
        if ch in "([{":
            stack.append(ch)
        else:
            if not stack or stack.pop() != pairs[ch]:
                return False
    return len(stack) == 0

print(is_valid("()[]{}"))    # True
print(is_valid("([)]"))      # False
print(is_valid("{[]}"))      # True
\end{lstlisting}

\section*{第四课：树——层级结构}

\subsection*{什么是树？}

树是一种\textbf{层级结构}的数据。你的文件系统就是一棵树：

\begin{verbatim}
C:\
├── Windows/
│   ├── System32/
│   └── ...
├── Program Files/
│   ├── Python/
│   └── ...
└── Users/
    └── ...
\end{verbatim}

\subsection*{二叉树}

二叉树是每个节点最多有两个子节点（左孩子和右孩子）的树。

\begin{lstlisting}[language=Python]
class TreeNode:
    def __init__(self, value):
        self.value = value
        self.left = None
        self.right = None

# 构建一棵树
root = TreeNode(1)
root.left = TreeNode(2)
root.right = TreeNode(3)
root.left.left = TreeNode(4)
root.left.right = TreeNode(5)
\end{lstlisting}

这棵树长这样：
\begin{verbatim}
    1
   / \
  2   3
 / \
4   5
\end{verbatim}

\subsection*{树的遍历}

\begin{lstlisting}[language=Python]
# 前序遍历：根 → 左 → 右
def preorder(node):
    if node:
        print(node.value, end=" ")
        preorder(node.left)
        preorder(node.right)

# 中序遍历：左 → 根 → 右
def inorder(node):
    if node:
        inorder(node.left)
        print(node.value, end=" ")
        inorder(node.right)

# 后序遍历：左 → 右 → 根
def postorder(node):
    if node:
        postorder(node.left)
        postorder(node.right)
        print(node.value, end=" ")

print("前序：", end="")
preorder(root)    # 1 2 4 5 3
print()
print("中序：", end="")
inorder(root)     # 4 2 5 1 3
print()
print("后序：", end="")
postorder(root)   # 4 5 2 3 1
\end{lstlisting}

\subsection*{二叉搜索树}

二叉搜索树的规则：对于每个节点，左子树的所有值都小于它，右子树的所有值都大于它。

\begin{lstlisting}[language=Python]
def insert(root, value):
    if not root:
        return TreeNode(value)
    if value < root.value:
        root.left = insert(root.left, value)
    else:
        root.right = insert(root.right, value)
    return root

def search(root, value):
    if not root or root.value == value:
        return root
    if value < root.value:
        return search(root.left, value)
    return search(root.right, value)
\end{lstlisting}

二叉搜索树的查找、插入、删除平均都是 $O(\log n)$——非常高效！

\section*{第五课：图——连接一切}

\subsection*{什么是图？}

图由\textbf{顶点}和\textbf{边}组成。社交网络（人是顶点，好友关系是边）、地图（地点是顶点，道路是边）、互联网都是图。

\subsection*{图的表示}

\begin{lstlisting}[language=Python]
# 用字典表示图（邻接表）
graph = {
    "A": ["B", "C"],
    "B": ["A", "D", "E"],
    "C": ["A", "F"],
    "D": ["B"],
    "E": ["B", "F"],
    "F": ["C", "E"]
}
\end{lstlisting}

\subsection*{广度优先搜索（BFS）}

从起点出发，先访问所有距离为 1 的节点，再访问距离为 2 的……就像水滴在纸上慢慢扩散。

\textbf{用途：} 最短路径、社交网络"你可能认识的人"。

\begin{lstlisting}[language=Python]
from collections import deque

def bfs(graph, start):
    visited = set()
    queue = deque([start])
    visited.add(start)
    
    while queue:
        node = queue.popleft()
        print(node, end=" ")
        
        for neighbor in graph[node]:
            if neighbor not in visited:
                visited.add(neighbor)
                queue.append(neighbor)

bfs(graph, "A")  # A B C D E F
\end{lstlisting}

\subsection*{深度优先搜索（DFS）}

从起点出发，沿着一条路走到黑，走不动了再退回来换条路。

\textbf{用途：} 迷宫求解、检测环、拓扑排序。

\begin{lstlisting}[language=Python]
def dfs(graph, start, visited=None):
    if visited is None:
        visited = set()
    visited.add(start)
    print(start, end=" ")
    
    for neighbor in graph[start]:
        if neighbor not in visited:
            dfs(graph, neighbor, visited)

dfs(graph, "A")  # A B D E F C
\end{lstlisting}

\section*{第六课：排序算法}

排序是最经典的算法问题——把一堆无序的数据排成有序的。

\subsection*{冒泡排序——最简单（也最慢）}

\begin{lstlisting}[language=Python]
def bubble_sort(arr):
    n = len(arr)
    for i in range(n):
        for j in range(n - 1 - i):
            if arr[j] > arr[j + 1]:
                arr[j], arr[j + 1] = arr[j + 1], arr[j]
    return arr

print(bubble_sort([64, 34, 25, 12, 22, 11, 90]))
\end{lstlisting}

思想：每轮把最大的数"冒"到最后。时间复杂度：$O(n^2)$。

\subsection*{选择排序——每次选最小的}

\begin{lstlisting}[language=Python]
def selection_sort(arr):
    n = len(arr)
    for i in range(n):
        min_idx = i
        for j in range(i + 1, n):
            if arr[j] < arr[min_idx]:
                min_idx = j
        arr[i], arr[min_idx] = arr[min_idx], arr[i]
    return arr
\end{lstlisting}

思想：每轮找出最小的元素放到最前面。时间复杂度：$O(n^2)$。

\subsection*{插入排序——像整理扑克牌}

\begin{lstlisting}[language=Python]
def insertion_sort(arr):
    for i in range(1, len(arr)):
        key = arr[i]
        j = i - 1
        while j >= 0 and arr[j] > key:
            arr[j + 1] = arr[j]
            j -= 1
        arr[j + 1] = key
    return arr
\end{lstlisting}

思想：把每个元素插入到已经排好序的部分里。对于接近有序的数据很快。

\subsection*{归并排序——分而治之}

\begin{lstlisting}[language=Python]
def merge_sort(arr):
    if len(arr) <= 1:
        return arr
    
    mid = len(arr) // 2
    left = merge_sort(arr[:mid])
    right = merge_sort(arr[mid:])
    
    return merge(left, right)

def merge(left, right):
    result = []
    i = j = 0
    while i < len(left) and j < len(right):
        if left[i] <= right[j]:
            result.append(left[i])
            i += 1
        else:
            result.append(right[j])
            j += 1
    result.extend(left[i:])
    result.extend(right[j:])
    return result
\end{lstlisting}

思想：把数组不断对半拆，拆到只剩一个元素，然后两两合并成有序的。时间复杂度：$O(n \log n)$。这是目前最好的通用排序算法之一。

\subsection*{快速排序——最快的通用排序}

\begin{lstlisting}[language=Python]
def quick_sort(arr):
    if len(arr) <= 1:
        return arr
    
    pivot = arr[len(arr) // 2]
    left = [x for x in arr if x < pivot]
    middle = [x for x in arr if x == pivot]
    right = [x for x in arr if x > pivot]
    
    return quick_sort(left) + middle + quick_sort(right)
\end{lstlisting}

思想：选一个"基准"，比基准小的放左边，比基准大的放右边，然后分别对左右做同样的事。平均 $O(n \log n)$，最坏 $O(n^2)$，但实际中通常是最快的。

\subsection*{排序算法对比}

\begin{center}
\begin{tabular}{|c|c|c|c|}
\hline
\textbf{算法} & \textbf{最好} & \textbf{平均} & \textbf{最坏} \\ \hline
冒泡排序 & $O(n)$ & $O(n^2)$ & $O(n^2)$ \\ \hline
选择排序 & $O(n^2)$ & $O(n^2)$ & $O(n^2)$ \\ \hline
插入排序 & $O(n)$ & $O(n^2)$ & $O(n^2)$ \\ \hline
归并排序 & $O(n \log n)$ & $O(n \log n)$ & $O(n \log n)$ \\ \hline
快速排序 & $O(n \log n)$ & $O(n \log n)$ & $O(n^2)$ \\ \hline
\end{tabular}
\end{center}

\section*{第七课：二分查找与二分答案}

\subsection*{二分查找}

你已经在上文中见过二分查找了。它只能在\textbf{有序}的数组中使用。

\begin{lstlisting}[language=Python]
def binary_search(arr, target):
    left, right = 0, len(arr) - 1
    
    while left <= right:
        mid = (left + right) // 2
        if arr[mid] == target:
            return mid
        elif arr[mid] < target:
            left = mid + 1
        else:
            right = mid - 1
    return -1

nums = [1, 3, 5, 7, 9, 11, 13]
print(binary_search(nums, 7))   # 3
print(binary_search(nums, 2))   # -1
\end{lstlisting}

每次把搜索范围缩小一半，所以是 $O(\log n)$。1000 个元素最多查 10 次，100 万个最多查 20 次。

\subsection*{练习：找平方根}

用二分法求一个整数的平方根（取整）。

\begin{lstlisting}[language=Python]
def sqrt_int(n):
    left, right = 0, n
    while left <= right:
        mid = (left + right) // 2
        if mid * mid <= n and (mid + 1) * (mid + 1) > n:
            return mid
        elif mid * mid < n:
            left = mid + 1
        else:
            right = mid - 1
    return 0

print(sqrt_int(16))   # 4
print(sqrt_int(20))   # 4 (4^2=16 <=20, 5^2=25>20)
\end{lstlisting}

\section*{第八课：递归——函数调自己}

递归是一个函数调用它自己。你已经在树的遍历和归并排序中见过它了。

\subsection*{递归三要素}

\begin{enumerate}
  \item \textbf{基准条件：} 什么时候停止递归（最简单的情况）
  \item \textbf{递归条件：} 把问题缩小一点（向基准条件靠近）
  \item \textbf{确保终止：} 每次递归都更接近基准条件
\end{enumerate}

\subsection*{例子：阶乘}

\begin{lstlisting}[language=Python]
def factorial(n):
    # 基准条件
    if n <= 1:
        return 1
    # 递归条件
    return n * factorial(n - 1)

print(factorial(5))  # 5*4*3*2*1 = 120
\end{lstlisting}

执行过程：
\begin{verbatim}
factorial(5) = 5 * factorial(4)
            = 5 * 4 * factorial(3)
            = 5 * 4 * 3 * factorial(2)
            = 5 * 4 * 3 * 2 * factorial(1)
            = 5 * 4 * 3 * 2 * 1
            = 120
\end{verbatim}

\subsection*{例子：斐波那契数列}

\begin{lstlisting}[language=Python]
# 朴素递归（很慢！O(2^n)）
def fib_slow(n):
    if n <= 1:
        return n
    return fib_slow(n - 1) + fib_slow(n - 2)

print(fib_slow(10))  # 55
# fib_slow(40) 会等很久！

# 带记忆的递归（快！O(n)）
def fib_fast(n, memo={}):
    if n in memo:
        return memo[n]
    if n <= 1:
        return n
    memo[n] = fib_fast(n - 1, memo) + fib_fast(n - 2, memo)
    return memo[n]

print(fib_fast(100))  # 354224848179261915075
\end{lstlisting}

\textbf{关键：} 朴素斐波那契做了大量重复计算。用"记忆化"（把算过的结果存起来）可以大大加快速度。这就是动态规划的雏形。

\section*{第九课：动态规划入门}

动态规划（Dynamic Programming，简称 DP）的核心思想：\textbf{把大问题拆成小问题，先解决小问题，把小问题的答案存起来，再用它们解决大问题。}

\subsection*{爬楼梯问题}

你爬楼梯，每次可以走 1 级或 2 级。问：爬到第 $n$ 级有多少种不同的走法？

\begin{itemize}
  \item 到第 1 级：1 种（走 1 步）
  \item 到第 2 级：2 种（1+1 或 2）
  \item 到第 3 级：可以从第 1 级走 2 步，或从第 2 级走 1 步，所以 = 1 + 2 = 3 种
  \item 到第 $n$ 级：$f(n) = f(n-1) + f(n-2)$ —— 又是斐波那契！
\end{itemize}

\begin{lstlisting}[language=Python]
def climb_stairs(n):
    if n <= 2:
        return n
    dp = [0] * (n + 1)
    dp[1] = 1
    dp[2] = 2
    for i in range(3, n + 1):
        dp[i] = dp[i - 1] + dp[i - 2]
    return dp[n]

print(climb_stairs(10))  # 89
\end{lstlisting}

这里 \texttt{dp[i]} 表示"爬到第 i 级有多少种方法"。我们从小到大把每个结果都算出来存好——这就是动态规划。

\subsection*{0/1 背包问题}

你有一个能装 10kg 的背包和三件物品，每件物品有重量和价值：

\begin{center}
\begin{tabular}{|c|c|c|}
\hline
\textbf{物品} & \textbf{重量} & \textbf{价值} \\ \hline
书 & 3kg & 50 元 \\ \hline
水 & 2kg & 30 元 \\ \hline
食物 & 5kg & 80 元 \\ \hline
\end{tabular}
\end{center}

怎么装价值最高？这是一个经典 DP 问题。

\begin{lstlisting}[language=Python]
def knapsack(weights, values, capacity):
    n = len(weights)
    dp = [[0] * (capacity + 1) for _ in range(n + 1)]
    
    for i in range(1, n + 1):
        for w in range(1, capacity + 1):
            if weights[i-1] <= w:
                dp[i][w] = max(dp[i-1][w], 
                              dp[i-1][w-weights[i-1]] + values[i-1])
            else:
                dp[i][w] = dp[i-1][w]
    
    return dp[n][capacity]

weights = [3, 2, 5]
values = [50, 30, 80]
print(knapsack(weights, values, 10))  # 160（书+水+食物全部装下）
\end{lstlisting}

\section*{总结：知识图谱}

\begin{verbatim}
数据结构与算法
├── 时间复杂度（大O）
├── 数据结构
│   ├── 数组（Python列表）
│   ├── 链表
│   ├── 栈（后进先出）
│   ├── 队列（先进先出）
│   ├── 树
│   │   ├── 二叉树
│   │   ├── 遍历（前/中/后序）
│   │   └── 二叉搜索树
│   └── 图
│       ├── 邻接表
│       ├── BFS
│       └── DFS
├── 算法
│   ├── 排序
│   │   ├── 冒泡排序 O(n²)
│   │   ├── 选择排序 O(n²)
│   │   ├── 插入排序 O(n²)
│   │   ├── 归并排序 O(n log n)
│   │   └── 快速排序 O(n log n)
│   ├── 二分查找 O(log n)
│   ├── 递归
│   └── 动态规划
│       ├── 记忆化
│       ├── 爬楼梯
│       └── 背包问题
\end{verbatim}

\textbf{下一步：} 数据结构与算法是编程的"内功"。掌握了这些，你可以在以下方向继续深入：
\begin{itemize}
  \item \textbf{信息学奥赛：} 更复杂的 DP、图论、数论
  \item \textbf{Python → C++：} 用 C++ 重新实现这些算法，感受性能差异
  \item \textbf{刷题平台：} LeetCode、Codeforces、洛谷
\end{itemize}

\section*{综合练习与答案}

\subsection*{基础题}

\begin{enumerate}
  \item 用栈判断括号字符串 \texttt{"\{\{()[]\}\}()"} 是否合法。
  \item 给一个无序列表 \texttt{[3, 7, 1, 9, 4, 2, 8, 5, 6]}，分别用手算和写程序验证冒泡排序和快速排序的过程。
  \item 用二分查找在 \texttt{[1, 3, 5, 7, 9, 11, 13, 15, 17, 19]} 中查找 13，写出每一步的 \texttt{left, right, mid} 值。
  \item 用递归实现 \texttt{sum\_list(lst)} 求列表所有元素的和（不能用 \texttt{sum()}）。
  \item 爬楼梯问题的变种：每次可以走 1 级、2 级或 3 级，求到第 $n$ 级的方法数。
\end{enumerate}

\subsection*{参考答案}

\textbf{练习 1：}
\begin{lstlisting}[language=Python]
def is_valid(s):
    stack = []
    pairs = {")": "(", "]": "[", "}": "{"}
    for ch in s:
        if ch in "([{":
            stack.append(ch)
        else:
            if not stack or stack.pop() != pairs[ch]:
                return False
    return len(stack) == 0

print(is_valid("{{()[]}}()"))  # True
\end{lstlisting}

\textbf{练习 3（二分查找过程）：}
\begin{verbatim}
初始：left=0, right=9, mid=4 → arr[4]=9 < 13 → left=5
第2步：left=5, right=9, mid=7 → arr[7]=15 > 13 → right=6
第3步：left=5, right=6, mid=5 → arr[5]=11 < 13 → left=6
第4步：left=6, right=6, mid=6 → arr[6]=13 ✓ 找到！
\end{verbatim}

\textbf{练习 4：}
\begin{lstlisting}[language=Python]
def sum_list(lst):
    if not lst:
        return 0
    return lst[0] + sum_list(lst[1:])

print(sum_list([1, 2, 3, 4, 5]))  # 15
\end{lstlisting}

\textbf{练习 5：}
\begin{lstlisting}[language=Python]
def climb_stairs_3(n):
    if n <= 1:
        return 1
    if n == 2:
        return 2
    dp = [0] * (n + 1)
    dp[0] = 1
    dp[1] = 1
    dp[2] = 2
    for i in range(3, n + 1):
        dp[i] = dp[i-1] + dp[i-2] + dp[i-3]
    return dp[n]

print(climb_stairs_3(10))  # 274
\end{lstlisting}

\section*{扩展项目}

\subsection*{项目一：表达式计算器}

实现一个能计算算术表达式的程序，如 \texttt{"(3 + 5) * 2 - 8 / 4"}。

\textbf{要求：}
\begin{enumerate}
  \item 支持加减乘除和括号
  \item 用"调度场算法"将中缀表达式转后缀表达式
  \item 用栈计算后缀表达式
\end{enumerate}

\textbf{关联模块：} 本项目综合运用了栈（本模块第三课）和递归思想。你在 \texttt{python/} 模块第一课学的字符串处理和本模块的栈可以完美配合。

\subsection*{项目二：文件搜索工具}

实现一个 \texttt{find\_in\_files(folder, keyword)} 函数，递归搜索文件夹中所有包含关键字的文件。

\textbf{关联模块：} 树的遍历（本模块第四课）+ 文件 IO（\texttt{python/} 模块第十课）+ 递归（本模块第八课）。你能在 \texttt{python/} 模块的记账本项目基础上做一层新的抽象。

\subsection*{项目三：排序算法性能对比}

写一个程序，生成 1000/5000/10000/50000 个随机数，分别用冒泡、选择、插入、归并、快排和 Python 内置的 \texttt{sorted()} 排序，记录每种算法的运行时间。绘图展示结果。

\textbf{关联模块：} 本模块的大 O 复杂度（第一课）+ \texttt{computer\_maths/} 模块的取模运算（生成随机数时用到）。实际运行结果应该和大 O 理论一致——数学预测和实验数据互相验证。

\section*{本模块与其他模块的关联}

\begin{itemize}
  \item \textbf{Python 基础} → 本模块：你在 Python 模块学到的列表、字典、函数是 ds-algo 的\textbf{前提条件}。链表（本模块第二课）是用 Python 的类实现的，如果你对类不熟悉，回去看 Python 模块的第八课。
  \item \textbf{计算机数学} ↔ 本模块：本模块的大 O 复杂度分析需要 \texttt{computer\_maths/} 的数列与求和知识；图的部分直接对应 \texttt{computer\_maths/} 的图论章节。两模块可以对照学习。
  \item 本模块 → \textbf{AI 入门}：AI 模块中神经网络的"全连接层"本质就是矩阵乘法（\texttt{computer\_maths/}）+ 映射函数。模型训练涉及梯度下降（优化算法），这和本模块的算法优化思想一脉相承。
  \item 本模块 → \textbf{Python → C++}：你把本模块的排序、搜索、栈、队列等算法用 C++ 重新实现一遍，就是最好的 C++ 练习。这也是 \texttt{python-to-cpp/} 模块的综合项目建议。
\end{itemize}

\end{document}
